%% ============================================================
%%  Capítulo 1 – Introducción
%% ============================================================
\chapter{Introducción}
\label{chap:introduccion}

\section{Propósito del Documento}
\label{sec:intro_proposito}

Este documento constituye la consolidación técnica definitiva del hardware de ALMA correspondiente al Entregable 3 (E3) del proyecto. Su propósito es responder a la pregunta central de esta etapa: \textbf{¿Cómo queda técnicamente definido ALMA antes del diseño físico final?}

El E3 marca el cierre formal de la fase de definición de arquitectura y selección de componentes, y da paso a la fase de enrutamiento físico (Entregable 4). A diferencia de los entregables anteriores, que establecían requisitos y especificaciones funcionales, este documento presenta decisiones de ingeniería concretas: componentes seleccionados con número de parte real, esquemáticos completados al 100\,\%, restricciones de fabricación configuradas en Altium Designer y reglas de integridad de señal listas para su aplicación en el layout.

El diseño aquí presentado es técnicamente defendible, ha pasado la verificación ERC (Electrical Rule Check) sin errores críticos, y está optimizado para la manufactura (Design for Manufacturing, DFM) con el proceso de fabricación del fabricante seleccionado.

\section{Relación con los Entregables Anteriores}
\label{sec:intro_relacion}

El Entregable 1 (E1) estableció los \textbf{requisitos funcionales, no funcionales, de interfaz y de restricción} del sistema ALMA. El Entregable 2 (E2) transformó esos requisitos en especificaciones técnicas cuantitativas y abrió cinco decisiones de diseño (D-01 a D-05) que no podían resolverse sin un análisis arquitectónico más profundo. El presente E3 cierra las cinco decisiones abiertas:

\begin{itemize}
  \item \textbf{D-01} -- Plataforma de cómputo: Amlogic A311D2 en implementación Chip-Down (\S\ref{sec:arq_soc}).
  \item \textbf{D-02} -- Tipo de micrófonos: MEMS digitales con interfaz PDM (\S\ref{sec:arq_audio}).
  \item \textbf{D-03} -- Estrategia de integración: Chip-Down vs. SoM (\S\ref{sec:arq_integracion}).
  \item \textbf{D-04} -- Conectividad y protocolo de domótica: WiFi + cliente IP (\S\ref{sec:arq_conectividad}).
  \item \textbf{D-05} -- Elemento de seguridad criptográfica: TrustZone integrado del A311D2 (\S\ref{sec:arq_seguridad}).
\end{itemize}

\section{Alcance y Exclusiones}
\label{sec:intro_alcance}

\textbf{El presente documento incluye:}
\begin{itemize}
  \item Esquemáticos completados al 100\,\% para la PCB-CTRL y la PCB-PWR.
  \item Lista de materiales (BOM) cerrada con números de parte concretos y precios de referencia.
  \item Árbol de potencia detallado con componentes generadores y cargas por riel.
  \item Restricciones del fabricante de PCB traducidas a reglas DRC en Altium Designer.
  \item Requisitos de integridad de señal (SI) e integridad de potencia (PI) para cada interfaz crítica.
  \item Arquitectura lógica del software y flujos de operación clave vinculados al hardware.
  \item Trazabilidad completa de los requisitos del E1 hacia las decisiones de diseño.
  \item Análisis de costos por volumen de producción.
  \item Registro de riesgos actualizado y decisiones sensibles para el enrutamiento.
\end{itemize}

\textbf{Quedan fuera del alcance de este entregable:}
\begin{itemize}
  \item El layout final de las PCBs (enrutamiento), que corresponde al Entregable 4.
  \item Los archivos de fabricación (Gerber, ODB++, IPC-2581).
  \item Las pruebas eléctricas sobre prototipo físico.
  \item La implementación del stack de software o firmware.
\end{itemize}

\section{Organización del Documento}
\label{sec:intro_organizacion}

El documento está estructurado en trece capítulos y cinco anexos. Los Capítulos~\ref{chap:arquitectura} y~\ref{chap:particion} describen la arquitectura final del sistema y la justificación de la partición en dos PCBs. Los Capítulos~\ref{chap:pcb_ctrl} y~\ref{chap:pcb_pwr} detallan cada tarjeta por separado. Los Capítulos~\ref{chap:fabricante} y~\ref{chap:si} documentan las restricciones de manufactura e integridad de señal. El Capítulo~\ref{chap:hw_sw} desarrolla la arquitectura lógica del software. Los Capítulos~\ref{chap:bom} a~\ref{chap:riesgos} cubren la lista de materiales, análisis de costos, trazabilidad de requisitos y gestión de riesgos. El Capítulo~\ref{chap:conclusiones} sintetiza el estado del proyecto y los puntos críticos para el enrutamiento. Los Anexos A--E contienen el BOM detallado, los cálculos de SI, el árbol de potencia completo y los esquemáticos exportados de Altium Designer.
